\chapter{Финальная заморозка статуса Version III.H}

После внутреннего и двух внешних red-team циклов фиксируется окончательный
статус третьего тома. Заморозка относится к минимальному hidden-sector
parent action и не переносится автоматически на исходную observed-world
программу Томов I--II.

\section{Определение замороженной модели}

Version III.H задаётся:
\begin{enumerate}
 \item finite algebra \(A_F=\mathbb C\oplus\mathbb C\);
 \item минимальным reality-замкнутым first-order bimodule square;
 \item KO6 completion с Pfaffian physical fermion count;
 \item faithful gauge quotient \(U(1)_{\mathrm{rel}}\);
 \item двумя opposite-charge scalar components \(x,z\);
 \item flattening potential
 \(\Tr(\Phi^2-\chi^2\mathbf1)^2\);
 \item full bosonic spectral trace \(\Tr_{\mathrm{spec}}\);
 \item одним dimensionful calibration scale \(\chi\).
\end{enumerate}

После canonical normalization
\begin{equation}
 \kappa=2,
 \qquad
 g^2(\mu_{\mathrm{spec}})=\frac38.
\end{equation}
Vacuum orbit имеет
\begin{equation}
 |x|=|z|=\frac{\chi}{\sqrt2},
 \qquad
 \arg(xz)=\pm\frac{\pi}{2}.
\end{equation}

\section{Доказанные результаты}

\begin{enumerate}
 \item first-order condition требует четырёхсекторный reality-square;
 \item finite potential выбирает equal moduli и maximal CP magnitude;
 \item scalar spectrum содержит три modes с
 \(m_s^2=4\chi^2\);
 \item две physical Dirac pairs имеют \(m_f=\chi\);
 \item hidden vector имеет \(m_A^2=3\chi^2\);
 \item linear и cubic \(U(1)\) anomalies сокращаются;
 \item physical zero-mode supertrace numerator равен
 \begin{equation}
  N_0=67;
 \end{equation}
 \item Coleman--Weinberg seed положителен:
 \begin{equation}
  B_0=\frac{67}{64\pi^2}>0;
 \end{equation}
 \item fermionic half-count выводится из Pfaffian;
 \item full bosonic trace сохраняет все normalized ratios;
 \item minimal direct-sum completion exactly decoupled от observed sector.
\end{enumerate}

\begin{theorem}[Mathematical closure of Version III.H]
Перечисленные algebra, representation, action, measure dictionary и
vacuum equations определяют согласованную anomaly-free
four-dimensional hidden-\(U(1)\) effective field theory. Внутренних
математических противоречий в замороженной версии не обнаружено.
\end{theorem}

\section{Допустимые входы и дискретные ветви}

Модель не является parameter-free в абсолютных единицах. Разрешён ровно
один dimensionful train datum:
\begin{equation}
 m_{\mathrm{ref}}=\chi.
\end{equation}
Эквивалентно он фиксирует spectral invariant
\begin{equation}
 \zeta_{\mathrm{mom}}
 =
 \frac{f_2}{f_0}(\Lambda R)^2
 =
 (\chi R)^2+\frac12.
\end{equation}

Выбор representative spectral profile внутри moment--cutoff rescaling
orbit не является дополнительным observable parameter. Он должен быть
зафиксирован как convention.

Две CP-сопряжённые ветви
\begin{equation}
 \phi=\pm\frac{\pi}{2}
\end{equation}
являются точным spontaneous degeneracy. Отсутствие orientation selector
не нарушает mathematical consistency: frozen theory содержит два
вырожденных vacuum sectors. Уникальный CP-sign не заявляется.

\section{Отрицательно закрытые ветви}

\begin{enumerate}
 \item flat internal-\(K\) spinorial lift не имеет fermion zero-modes;
 \item direct Standard Model readout minimal algebra невозможен;
 \item charged-lepton identification двух Dirac pairs проваливается из-за
 exact degeneracy;
 \item neutrino mass operator и rank selector отсутствуют;
 \item scalar, kinetic-mixing и neutrino portals не выводятся;
 \item minimal direct-sum portal coefficients точно равны нулю;
 \item \(S_{\mathrm{vac}}\) Тома II не выведен из hidden functional;
 \item minimal EW/QCD repair routes Тома II остаются отрицательными;
 \item natural spectral-function menu не выбирает absolute scale.
\end{enumerate}

\section{Что не является блокером Version III.H}

Полный bosonic spectral action использует
\(\Tr_{\mathrm{spec}}=\Tr_{H_8}\). Поэтому universal bosonic
orbit half-trace больше не является assumption или незакрытым condition.
Pfaffian coefficient \(1/2\) применяется только к physical fermion
measure.

Общий множитель bosonic action зависит от spectral normalization, но
relative masses и couplings invariant. После одного train scale это
обычная EFT calibration, а не скрытая подгонка.

\section{Физический статус}

\begin{theorem}[Phenomenological status]
Version III.H не является phenomenologically testable model наблюдаемого
мира в текущем виде. Exact nongravitational decoupling означает отсутствие
laboratory portal, а gravitational normalization не выведена.
\end{theorem}

Следовательно, корректная формулировка:
\begin{equation}
 \boxed{
 \begin{gathered}
 \text{математически замкнутая one-scale hidden-sector EFT},\\
 \text{не unified theory наблюдаемого микромира},\\
 \text{не модель с текущими laboratory predictions}.
 \end{gathered}}
\end{equation}

\section{Связь с Томами I и II}

Том I сохраняется как философско-методологическая программа. Том II
сохраняет geometric ядро, conditional sector models и отрицательные
falsification results. Том III.H является новой конкретной реализацией
части идеи structural consistency, но не common source всех прежних
формул.

Не построены:
\begin{equation}
 \Gamma_{\mathrm{unif}}
 \longrightarrow
 \left\{
 S_{\mathrm{vac}},
 \Gamma_{\mathrm{hidden}},
 \Gamma_{\mathrm{EW/QCD}},
 \Gamma_\nu
 \right\}.
\end{equation}
Поэтому cross-tome unification остаётся задачей отдельной версии.

\section{Условия переоткрытия}

Hidden branch может быть расширен только если заранее выведено одно из:
\begin{enumerate}
 \item connector bimodule, фиксирующее portal coefficient;
 \item common algebra, порождающая bi-charged states;
 \item orientation-odd invariant, выбирающий CP branch;
 \item gravitational measure с вычисленным Planck normalization.
\end{enumerate}

Observed-sector reconstruction открывается только новой Version IV.SM,
которая должна независимо задать:
\begin{enumerate}
 \item \(SU(3)\times SU(2)\times U(1)_Y\) representation;
 \item anomaly-free generation structure;
 \item neutrino mass graph;
 \item common functional bridge к electromagnetic sector;
 \item минимум два preregistered physical blind tests.
\end{enumerate}

\begin{protocol}[Final freeze]
С 10 августа 2026 года Version III.H заморожена. Запрещено:
\begin{enumerate}
 \item добавлять portal coefficients без новой representation;
 \item переименовывать hidden modes в Standard Model particles post hoc;
 \item использовать \(S_{\mathrm{vac}}\), число \(23\) или EW/QCD formulas
 как доказательство hidden action;
 \item считать один train scale blind prediction.
\end{enumerate}
Все дальнейшие observed-world расширения относятся к Version IV.SM.
\end{protocol}

Машинный freeze ledger сохранён в
\path{s2t_v3_final_status_freeze_results.json}.